\section{Repaso de los objetos centrales}

\begin{importantbox}{Función de valor, función Q y función de ventaja}
\begin{itemize}
    \item \textbf{Función de valor} $V^\pi(s)$: la recompensa futura esperada al partir del estado $s$ y seguir $\pi$.
    \item \textbf{Función Q} $Q^\pi(s, a)$: la recompensa futura esperada al ejecutar $a$ desde $s$ y luego seguir $\pi$.
    \item \textbf{Función de ventaja} $A^\pi(s, a) = Q^\pi(s, a) - V^\pi(s)$: cuánto mejor es la acción $a$ en comparación con el desempeño promedio de la política $\pi$.
\end{itemize}
Relación clave: $V^\pi(s) = \mathbb{E}_{a \sim \pi(\cdot|s)}\left[Q^\pi(s, a)\right]$
\end{importantbox}

\subsection{Resumen de la sección}

Las tres funciones $V$, $Q$ y $A$ describen desde distintos ángulos qué tan buena es una política. La función de ventaja $A$ resulta especialmente útil porque mide directamente qué tan buena es una acción en comparación con el desempeño promedio de la política.
